\section{The dummy neighborhood meets only the initial layer}
\label{sec:area_law_dummy_contacts}

The dummy tile is the initial neighborhood $N_{k_0}$. Its distance
bound distinguishes it from the later layers; see
\cite[Section~11]{OpenAI2026AreaLaw}. The conclusions below concern
closures, with distances measured in the sup metric.

% Source: 10-geometry.tex, lines 160--191 and 299--306.

\begin{theorem}[Distance from the neighborhood to an endpoint]
    \label{thm:area_law_neighborhood_endpoint_distance}
    \lean{TNLean.PEPS.AreaLaw.Geometry.dyadicNeighborhood_exists_dist_le}
    \leanok
    \uses{def:area_law_dyadic_layers}
    For arbitrary origin, finite endpoint set $Z$, nonnegative integer
    width $C_0$ and index $k$, every $x\in\overline{N_k}$ admits
    an endpoint $z\in Z$ with
    \begin{align}
        d_\infty(x,z)\le(C_0+1)2^k.
    \end{align}
\end{theorem}
\begin{proof}\leanok
    \uses{thm:area_law_dyadic_closures,thm:area_law_dyadic_indices,
        thm:area_law_dyadic_cell_distance}
    The point lies in a closed scale-$k$ cell whose index is within
    $C_0$ of an occupied index in each coordinate. Choose an endpoint
    in that occupied cell and apply the closed-cell distance bound.
\end{proof}

\begin{theorem}[Separation from subsequent layers]
    \label{thm:area_law_dummy_later_layer_separation}
    \lean{TNLean.PEPS.AreaLaw.Geometry.dyadicNeighborhood_dist_later_layer}
    \leanok
    \uses{def:area_law_dyadic_layers}
    Under the same hypotheses, let $k\ge k_0+1$,
    $x\in\overline{N_{k_0}}$, and $y\in\overline{D_k}$. Then
    \begin{align}
        (C_0-1)2^{k_0}\le d_\infty(x,y).
    \end{align}
    In particular, this is a positive separation when $C_0\ge2$.
\end{theorem}
\begin{proof}\leanok
    \uses{thm:area_law_neighborhood_endpoint_distance,
        thm:area_law_dyadic_layer_distance}
    Choose $z\in Z$ with
    $d_\infty(x,z)\le(C_0+1)2^{k_0}$. The layer bound gives
    $C_0 2^k\le d_\infty(y,z)$, while $2^k\ge2\cdot2^{k_0}$.
    The triangle inequality now yields
    \begin{align}
        d_\infty(x,y) & \ge C_0 2^k-(C_0+1)2^{k_0}
        \ge(C_0-1)2^{k_0}.
    \end{align}
\end{proof}

\begin{theorem}[Only the initial layer can meet the dummy]
    \label{thm:area_law_dummy_contact_index}
    \lean{TNLean.PEPS.AreaLaw.Geometry.dyadicNeighborhood_contact_layer_eq}
    \leanok
    \uses{def:area_law_dyadic_layers}
    Suppose that $C_0\ge2$ and $k\ge k_0$. If
    $\overline{N_{k_0}}\cap\overline{D_k}$ is nonempty, then $k=k_0$.
\end{theorem}
\begin{proof}\leanok
    \uses{thm:area_law_dummy_later_layer_separation}
    If $k\ne k_0$, then $k\ge k_0+1$. A shared point has distance
    zero from itself, contradicting the positive separation.
\end{proof}
